\documentclass[a4paper,12pt]{letter}

\usepackage{amsmath}
\usepackage{amssymb}
\usepackage{nopageno}
\usepackage{amsmath}
\usepackage{enumitem}
% \usepackage{mathtools}
\usepackage{eqnarray,amsmath}
\usepackage[a4paper, total={7.5in, 11.25in}]{geometry}
%\usepackage[margin=2.5cm]{geometry}
\usepackage[utf8]{inputenc}
\usepackage[T1]{fontenc}
\usepackage{lmodern}
%\usepackage{amsmath, amssymb}
\usepackage{setspace}
\setstretch{1.2}

\begin{document}

\begin{center}
(KTXFI2EBNF) Physics II. Homework No. 3 Solutions\\
November 25, 2025
\end{center}

\begin{enumerate}

\item We study an electron in a radioactive particle beam as it passes through the laboratory. It is found that, as measured in the laboratory, the average lifetime of the particle is $25\,ns$; after this time the particle decays. The average lifetime measured in the reference frame attached to the particles is $7.5\,ns$. What is the propagation velocity of the particle beam?\\

The Lorentz factor is
$$ \gamma = \frac{\Delta t}{\Delta \tau} = \frac{25}{7.5} = \frac{10}{3} \approx 3.3333333. $$
Using $\gamma = 1/\sqrt{1 - v^2/c^2}$ we obtain
$$ v = c\sqrt{1 - \frac{1}{\gamma^2}} = c\sqrt{1 - \left(\frac{3}{10}\right)^2} \approx 0.953939\,c. $$
Thus the beam propagates at approximately $0.95394$ times the speed of light.

\item A certain bacterium doubles its population in 25 days. Two such bacteria are placed in a spaceship and sent on an 800-day space journey. During the trip, the spaceship moves at a constant speed of $0.995\,c$. How many bacteria will be on board at the moment of return?\\

The proper time experienced on board is shorter than the laboratory time. The ship frame time is
$$ \Delta \tau = \Delta t\sqrt{1 - v^2/c^2} = 800\sqrt{1 - (0.995)^2}\ \text{days} \approx 79.89994\ \text{days}. $$
Exponential growth with doubling time 25 days gives
$$ N = N_0 \, 2^{\Delta \tau / 25} = 2 \cdot 2^{79.89994/25} \approx 18.33. $$
Approximately 18 bacteria will be on board at return (the continuous model gives about 18.3 individuals).

\item A particle moves with a velocity such that $\frac{v}{c} = 0.9999$. Determine the value of $\frac{m}{m_{0}}$.\\

The relativistic mass factor equals the Lorentz factor
$$ \frac{m}{m_0} = \gamma = \frac{1}{\sqrt{1 - v^2/c^2}} = \frac{1}{\sqrt{1 - (0.9999)^2}} \approx 70.712446. $$

\item Determine the relativistic mass and the velocity of an electron whose kinetic energy is 1 MeV ($ 1.6 \cdot 10^{-13} \, \text{J} $). The rest mass of the electron is $9.11 \cdot 10^{-31}$ kg.\\

From the relation between kinetic energy and gamma factor,
$$ E_{kin} = (\gamma - 1)m_0 c^2, $$ 
we obtain
$$ \gamma = 1 + \frac{E_{kin}}{m_0 c^2} \approx 1 + \frac{1.0\cdot 10^{6}\ \text{eV}\cdot 1.6\cdot 10^{-19}\ \text{J/eV}}{(9.11\cdot 10^{-31}\,\text{kg})(299792458\,\text{m/s})^2}
\approx 2.9541604. $$
The relativistic mass is
$$ m = \gamma m_0 \approx 2.9541604 \cdot 9.11\cdot 10^{-31}\ \text{kg} \approx 2.69124\cdot 10^{-30}\ \text{kg}. $$
The speed follows from
$$ v = c\sqrt{1 - \frac{1}{\gamma^2}} \approx 0.94096436\,c \approx 2.82094\cdot 10^{8}\ \text{m/s}. $$

\item Calculate the rest energy of a proton, corresponding to its rest mass of $1.67 \cdot 10^{-27}$ kg, in joules and in electronvolts. ($1\,\text{MeV} = 1.6 \cdot 10^{-13}\text{J}$)\\

The rest energy is
$$ E = m c^2 = (1.67\cdot 10^{-27})(299792458)^2 \approx 1.50092115\cdot 10^{-10}\ \text{J}. $$
In mega-electronvolts,
$$ E \approx \frac{1.50092115\cdot 10^{-10}}{1.6\cdot 10^{-13}}\ \text{MeV} \approx 938.08\ \text{MeV}. $$

\end{enumerate}

\bigskip

\rightline{Dr. Gabor FACSKO}
\rightline{\textit{facsko.gabor@uni-obuda.hu}}

\end{document}
